% BANKS-SOURCE: pdf=142 print=132 kind=solution-section id=chapter08-numbered-solutions
\begin{bankssolutions}{Solutions to Problems for Chapter 8}

% BANKS-SOURCE: pdf=142 print=132 kind=solution id=8.1
\BanksSolution{8.1}
Write the BRST differential as $s$, so that
$\delta_{\rm BRST}=\epsilon s$.  It is useful to combine the ghosts with the
generators into the odd Lie-algebra-valued field $c=c^aT_a$.  With a choice of
sign convention for the gauge transformation, the algebraic part of the BRST
rules is
\[
  sX=[X,c],\qquad sc=-\frac12[c,c]
\]
for an adjoint-valued ordinary field $X$.  For a field $\varphi$ in an
arbitrary representation the corresponding formula is $s\varphi=-c\varphi$.
The opposite gauge-parameter convention reverses the three gauge-action signs
together. In components these rules reproduce the transformation in the
question after the same convention is used for the generators and structure
constants.

For a matter field, the graded Leibniz rule gives
\begin{align*}
  s^2\varphi
    &=-(sc)\varphi+c(s\varphi)\\
    &=\frac12[c,c]\varphi-c^2\varphi=0,
\end{align*}
because $[c,c]=2c^2$ for an odd matrix $c$.  Equivalently, in components,
the first term is proportional to $f^a{}_{bc}c^bc^cT_a\varphi$, while the
second term reduces, using $c^bc^a=-c^ac^b$, to the same expression with the
opposite sign.  The same calculation applies to every tensor product or
direct-sum representation.

For the gauge potential one has
\[
  sA_\mu=D_\mu c=\partial_\mu c+[A_\mu,c].
\]
Using the graded derivation rule and
$D_\mu[c,c]=2[D_\mu c,c]$, one obtains
\[
  s^2A_\mu=D_\mu(sc)+[sA_\mu,c]
  =-\frac12D_\mu[c,c]+[D_\mu c,c]=0.
\]
The Jacobi identity is precisely what makes the same cancellation work in
$s^2c$.  Finally,
\[
  s\bar c=N,\qquad sN=0
\]
immediately imply $s^2\bar c=s^2N=0$.  Thus $s^2=0$ on all fields, and the
charge which implements it obeys $Q_{\rm BRST}^2=0$.

% BANKS-SOURCE: pdf=142 print=132 kind=solution id=8.2
\BanksSolution{8.2}
For this free Maxwell calculation we work in $D=4$ with
$\eta=\operatorname{diag}(-,+,+,+)$ and use
$p\cdot x=-p^0x^0+\mathbf p\cdot\mathbf x$. Positive-frequency modes carry
$\ee^{+\ii p\cdot x}$.
For the abelian theory the BRST rules can be chosen as
\[
  sA_\mu=\partial_\mu c,\qquad sc=0,\qquad s\bar c=N,
  \qquad sN=0.
\]
The gauge-fixed Lagrangian, before eliminating $N$, may be written as
\[
  \mathcal L=-\frac14F_{\mu\nu}F^{\mu\nu}
  +N\,\partial_\mu A^\mu+\frac\kappa2N^2
  +\partial_\mu\bar c\,\partial^\mu c,
\]
up to an overall sign in the ghost term that can be absorbed into $\bar c$.
Eliminating $N$ gives $N=-\partial\cdot A/\kappa$ and the usual term
$-(\partial\cdot A)^2/(2\kappa)$.  Noether's theorem gives the charge in a
form that does not depend on a choice of mode normalization,
\[
  Q_{\rm BRST}=\int\dd^{D-1}\mathbf x\,
  \left(\Pi_A^{\nu}\partial_\nu c+\Pi_{\bar c}N\right),
  \qquad
  \Pi_A^{\nu}=-F^{0\nu}+\eta^{0\nu}N,
  \quad \Pi_{\bar c}=\partial^0c .
\]
An integration by parts and the free equations give equivalent expressions
differing by a surface term.  In particular, the charge has the defining
commutators
\[
  [Q_{\rm BRST},A_\mu]=\ii\partial_\mu c,
  \qquad \{Q_{\rm BRST},\bar c\}=\ii N,
\]
with the factors of $\ii$ changed if the charge is defined with the opposite
Noether convention.

Let $k^0=|\mathbf k|$ and choose four polarization vectors, with $\epsilon_1$
and $\epsilon_2$ transverse to $\mathbf k$, $\epsilon_3$ longitudinal, and
$\epsilon_0$ time-like.  The free fields are expanded as
\[
  A_\mu(x)=\int\frac{\dd^{D-1}\mathbf k}{(2\pi)^{D-1}\sqrt{2k^0}}
  \sum_{\lambda=0}^3\epsilon_\mu{}^{(\lambda)}(k)
  \left[a_\lambda(k)\ee^{+\ii k\cdot x}
       +a_\lambda^\dagger(k)\ee^{-\ii k\cdot x}\right],
\]
where the oscillator metric is the Minkowski metric in polarization space.
The scalar ghosts have the same plane-wave normalization.  Denote their
annihilation operators by $c(k)$ and $b(k)$, with the conjugate pairing
$b=c^\dagger$ as stipulated in the question.  Substitution in the Noether
charge, followed by the free equations, gives, up to a nonzero normalization
and an overall sign,
\[
  Q_{\rm BRST}=\int\frac{\dd^{D-1}\mathbf k}{(2\pi)^{D-1}}
  \left[(a_0^\dagger(k)-a_3^\dagger(k))c(k)
       +c^\dagger(k)(a_0(k)-a_3(k))\right].
  \tag{8.2.1}
\]
The relative sign between $a_0$ and $a_3$ follows from
$k^\mu\epsilon_\mu^{(0,3)}$; changing the polarization convention merely
changes the displayed null combination.

Equation (8.2.1) exhibits the quartet mechanism.  The two transverse
oscillators commute with $Q_{\rm BRST}$.  The combination
$(a_0^\dagger-a_3^\dagger)|0\rangle$ is exact, since it is proportional to
$Q_{\rm BRST}c^\dagger|0\rangle$.  The complementary time-like/longitudinal
combination is paired with the ghost by the second term in (8.2.1), and the
anti-ghost completes the corresponding BRST doublet.  The same pairing applies
mode by mode and persists for multiparticle states by the graded Leibniz rule.
Consequently
\[
  H_{\rm BRST}=\frac{\ker Q_{\rm BRST}}
  {\operatorname{im}Q_{\rm BRST}}
  =\mathcal F(a_1^\dagger,a_2^\dagger)|0\rangle .
\]
The nontrivial cohomology is therefore generated exactly by the two transverse
photon polarizations.  States containing the time-like, longitudinal, ghost,
or anti-ghost excitations are either not closed or belong to a BRST-trivial
quartet.

% BANKS-SOURCE: pdf=143 print=133 kind=solution id=8.3
\BanksSolution{8.3}
Use $\eta=\operatorname{diag}(-,+,\ldots,+)$ and the positive-frequency
convention $\ee^{+\ii p\cdot x}$, so a massive on-shell momentum obeys
$p^2=-m_A^2$.
Let $f^a{}_i=(T^a v)_i$ and retain only the terms linear in the fluctuations.
The BRST transformations in an $R_\kappa$ gauge are
\[
  sA_\mu{}^a=\partial_\mu c^a,\qquad
  s\Delta_i=g_af_i{}^a c^a,\qquad
  sc^a=0,qquad s\bar c^a=N^a,qquad sN^a=0,
  \tag{8.3.1}
\]
where the sign of the second equation follows the convention for $D_\mu$.
The ghost and would-be Goldstone masses are $\kappa m_A^2$, while the
physical vector mass matrix is
\[
  (m_A^2)^{ab}=g_ag_bf_i{}^af_i{}^b .
\]

For a one-vector state with polarization $\epsilon_\mu$ and momentum $p$,
the first equation in (8.3.1) gives
\[
  s\bigl(\epsilon^\mu A_\mu{}^{a\dagger}(p)|0\rangle\bigr)
  \;\propto\;(\epsilon\cdot p)c^{a\dagger}(p)|0\rangle .
\]
An on-shell massive vector has $D-1$ polarizations satisfying
$p\cdot\epsilon=0$, including its physical longitudinal polarization. These
states are BRST closed. In each broken-generator sector, the polarization
parallel to $p_\mu$, the would-be Goldstone mode, the ghost, and the anti-ghost
form a contractible quartet. More explicitly, $Q_{\rm BRST}$ maps the
$p$-polarized vector and Goldstone span to the ghost, and maps the anti-ghost
back into that same bosonic span. Its image at ghost number zero is therefore
contained in the unphysical vector--Goldstone span and has zero projection on
the $D-1$ $p\cdot\epsilon=0$ vector states. Those states are closed and
non-exact. This argument also applies at $\kappa=1$, where the quartet masses
happen to coincide with the physical-vector mass.

For an unbroken generator, $m_A=0$.  The condition $p\cdot\epsilon=0$ leaves
$D-1$ closed polarizations, but the polarization proportional to $p_\mu$ is
BRST exact.  Quotienting by exact states leaves the $D-2$ transverse massless
polarizations.

The scalar transformation in (8.3.1) shows directly which scalar states
survive.  A fluctuation $u_i\Delta_i$ is closed precisely when
\[
  u_if_i{}^a=0\qquad\text{for every broken generator }a .
\]
These are the directions orthogonal to the would-be Goldstone subspace.  They
are the physical Higgs excitations.  The scalar combinations $f_i{}^a\Delta_i$
are paired with the broken-sector ghosts and the unphysical vector
polarizations; their $R_\kappa$ masses are equal, so they belong to BRST
quartets.  Thus the BRST cohomology contains $D-1$ polarizations of every
massive vector, $D-2$ transverse polarizations of every massless vector, and
exactly the scalar excitations orthogonal to all would-be Goldstone directions.

% BANKS-SOURCE: pdf=143 print=133 kind=solution id=8.4
\BanksSolution{8.4}
This Wilson-loop computation is four-dimensional and Euclidean.
Take the rectangle to have long sides at $\mathbf x=0$ and
$\mathbf x=\mathbf L$, with $|\mathbf L|=L$.  The Euclidean covariant-gauge
propagator is, in this fixed $D=4$ specialization,
\begin{equation}
\begin{aligned}
  D_{\mu\nu}{}^{ab}(q)
  &=\frac{\delta^{ab}}{q^2}
  \left[\delta_{\mu\nu}-(1-\kappa)\frac{q_\mu q_\nu}{q^2}\right],\\
  \langle A_\mu{}^a(x)A_\nu{}^b(y)\rangle
  % CONVENTION-SCOPE: D=4 reason=Euclidean Wilson-loop Fourier measure
  &=g^2\int\frac{\dd^4q}{(2\pi)^4}
  \ee^{\ii q\cdot(x-y)}D_{\mu\nu}{}^{ab}(q).
\end{aligned}
\tag{8.4.1}
\end{equation}
Expanding the path-ordered exponential to second order gives
\[
  \langle W_R(C)\rangle
  =\operatorname{Tr}_R\left[1-
  \frac{g^2}{2}\oint_C\dd x^\mu\oint_C\dd y^\nu\,
  T_R^aT_R^b\,D_{\mu\nu}{}^{ab}(x-y)+O(g^4)\right].
\]
The longitudinal term in (8.4.1) is a double total derivative along the
closed contour.  Its integral vanishes, or equivalently it vanishes by current
conservation.  The two short spatial sides give terms that remain finite as
$T\to\infty$.  The two long sides give a self-energy proportional to $T$ and
independent of $L$, together with the interaction between the two sources.
Here $D_{\mu\nu}{}^{ab}(x-y)$ denotes the Fourier transform of the first kernel
in (8.4.1).
In the same fixed $D=4$ specialization, using
\[
  \int\frac{\dd^3\mathbf q}{(2\pi)^3}
  \frac{\ee^{\ii\mathbf q\cdot\mathbf L}}{\mathbf q^2}
  =\frac{1}{4\pi L}\equiv\frac{c}{L},
\]
the interaction is the Coulomb operator
\[
  V_{\rm int}(L)=\frac{cg^2}{L}\,t_R^a t_{\bar R}^a .
  \tag{8.4.2}
\]
The second line runs backwards in Euclidean time.  Reversing its path changes
the ordered product to the conjugate representation, with
$t_{\bar R}^a=-(t_R^a)^T$.  This is the only representation-theoretic effect of
path ordering at this order; commutators generated by ordering first contribute
at higher powers of $g$.

For $T\gg L$, the two long lines propagate for a long Euclidean time.  Their
color indices are consequently an effective Hilbert space
$R\otimes\bar R$, and the single trace around the rectangle becomes
\[
  \langle W_R(T,L)\rangle
  =\operatorname{Tr}_{R\otimes\bar R}
  \exp\!\left[-T\frac{cg^2}{L}\,t_R^at_{\bar R}^a\right]
  \ee^{-T\infty}\,[1+O(L/T,g^4)].
  \tag{8.4.3}
\]
The last factor records the two $L$-independent source self-energies.  One can
obtain (8.4.3) directly by taking each long side to be an infinite Wilson
line, one future-directed and one past-directed, and then reversing the
past-directed line into $\bar R$.

Decompose the tensor product into irreducibles,
$R\otimes\bar R=\bigoplus_Qm_QQ$.  On the $Q$ summand,
\[
  (t_R^a+t_{\bar R}^a)^2=C_Q,
  \qquad
  t_R^at_{\bar R}^a
  =\frac12\bigl(C_Q-C_R-C_{\bar R}\bigr).
  \tag{8.4.4}
\]
Therefore
\[
  \langle W_R\rangle
  =\ee^{-T\infty}\sum_Qm_Q\dim Q
  \exp\!\left[-T\frac{cg^2}{2L}
  (C_Q-C_R-C_{\bar R})\right].
\]
Since $C_Q\geq0$, the smallest potential is attained for the singlet
$Q=\mathbf1$, which is present in $R\otimes\bar R$.  It has
$t_R^at_{\bar R}^a=-C_R$ and hence the attractive potential
$-cg^2C_R/L$.  For a baryon in SU(3),
$3\otimes3=\bar3\oplus6$; choosing the $\bar3$ channel for two quarks and
then combining it with the third $3$ gives
$\bar3\otimes3=\mathbf1\oplus8$.  The singlet channel is the most attractive,
which yields the large-angular-momentum picture described in the question.

% BANKS-SOURCE: pdf=143 print=133 kind=solution id=8.5
\BanksSolution{8.5}
Let
\[
  d_R^{abc}=\operatorname{tr}_R\bigl(T^a\{T^b,T^c\}\bigr).
\]
The conjugate representation has generators
$T_{\bar R}^a=-(T_R^a)^*$, so Hermiticity gives
\begin{align*}
  d_{\bar R}^{abc}
    &=\operatorname{tr}\left((-T_R^a)^*
      \{(-T_R^b)^*,(-T_R^c)^*\}\right)\\
    &=-\left[\operatorname{tr}_R
      \bigl(T_R^a\{T_R^b,T_R^c\}\bigr)\right]^*
     =-d_R^{abc}.
\end{align*}
The last equality uses the fact that the trace of a Hermitian matrix times an
anticommutator of Hermitian matrices is real.  If $R$ is real or pseudo-real,
there is a unitary intertwiner $U$ with
$T_{\bar R}^a=U T_R^aU^\dagger$.  The trace is invariant under this unitary
equivalence, so $d_{\bar R}^{abc}=d_R^{abc}$.  Combining the two relations
gives $d_R^{abc}=-d_R^{abc}=0$.  The argument uses only compactness, which
ensures a unitary realization of the representation, and applies equally to
real and pseudo-real representations.

% BANKS-SOURCE: pdf=143 print=133 kind=solution id=8.6
\BanksSolution{8.6}
The Wess--Zumino boundary density in this solution is four-dimensional; its
Chern--Simons extension is five-dimensional.
Let the first $M$ left-handed fermions transform as a fundamental of
$SU_L(M)$ and have U(1) charge $+1$.  The other $M$ left-handed fields are
the charge conjugates of right-handed fermions: they have U(1) charge $-1$
and transform as an anti-fundamental of $SU_R(M)$.  This is the chiral form of
the usual left--right flavor symmetry.  Write their background connection in
block form as
\[
  \mathcal A=\begin{pmatrix}
    A_1\mathbf 1_M+A_L&0\\[2pt]
    0&-A_1\mathbf 1_M-A_R^{T}
  \end{pmatrix},
  \qquad \operatorname{tr}A_L=\operatorname{tr}A_R=0 .
  \tag{8.6.1}
\]
The transpose in the second block is the anti-fundamental action.

Restricting to the diagonal subgroup sets $A_L=A_R=A_D$.  The two chiral
blocks then contribute opposite cubic traces,
\[
  \operatorname{tr}_M(A_1+A_D)^3
  +\operatorname{tr}_M(-A_1-A_D^T)^3=0 .
  \tag{8.6.2}
\]
The same cancellation applies to the U(1)$^3$, U(1)$SU_D(M)^2$, and mixed
gravitational anomalies.  Thus $H=\mathrm{U}(1)\times SU_D(M)$ is
anomaly-free, whereas the axial combinations of the two SU($M$) factors and
the U(1) direction orthogonal to the vector combination can have anomalies.

Here is an explicit descent representative for the Wess--Zumino density.
It also fixes the normalization conventions.  Use anti-Hermitian connection
forms, $F=\dd A+A^2$, and define
\[
  \omega_5(A)=\operatorname{tr}_{\mathcal R}
  \left(A(\dd A)^2+\frac32A^3\dd A+\frac35A^5\right),
  \qquad \dd\omega_5(A)=\operatorname{tr}_{\mathcal R}F^3,
  \tag{8.6.3}
\]
where the trace is over the $2M$ chiral fermions.  The overall coefficient is
the one multiplying the fermion anomaly; with the standard one-Weyl
normalization it is $\ii/(24\pi^2)$, while in the normalization of the
question it is fixed by the prefactor chosen for $S_{\rm CS}$.  Under
\[
  A^\Omega=\Omega^{-1}A\Omega+\Omega^{-1}\dd\Omega,
  \qquad \vartheta=\dd\Omega\,\Omega^{-1},
\]
the finite descent formula is
\[
  \omega_5(A^\Omega)-\omega_5(A)
  =\dd\alpha_4(A,\Omega)
   +\frac1{10}\operatorname{tr}_{\mathcal R}(\Omega^{-1}\dd\Omega)^5,
  \tag{8.6.4}
\]
with
\[
  \alpha_4(A,\Omega)=
  -\frac12\operatorname{tr}_{\mathcal R}\!\left[
    \vartheta\bigl(A\dd A+\dd A\,A+A^3\bigr)
    +\vartheta A\vartheta A+2\vartheta^3A\right].
  \tag{8.6.5}
\]
All products in (8.6.3)--(8.6.5) are wedge products.  Hence the local
four-dimensional term in the variation of the Chern--Simons action is
\[
  \mathcal L_{\rm WZ}(A,\Omega)=\mathcal C_{\mathcal R}\,\alpha_4(A,\Omega),
  \tag{8.6.6}
\]
where $\mathcal C_{\mathcal R}$ is the Chern--Simons prefactor.  The last
term in (8.6.4) is the usual winding-number term.  It is independent of $A$,
is defined modulo the quantized Chern--Simons period, and is part of the
global Wess--Zumino functional when a five-dimensional extension is needed.
Equations (8.6.5)--(8.6.6) are local in the boundary gauge field, $\Omega$,
and their first derivatives, as required.

The descent forms obey the group-composition identity
\[
  \alpha_4(A,\Omega h)=\alpha_4(A,\Omega)
  +\alpha_4(A^\Omega,h)+\dd(\cdots).
  \tag{8.6.7}
\]
For $h\in H$, the second term vanishes because every cubic invariant of the
anomaly-free subgroup vanishes by (8.6.2).  The possible exact term integrates
to a boundary-of-boundary contribution.  Thus $\mathcal L_{\rm WZ}$ depends
only on the right coset $[\Omega]\in G_F/H$.

Finally let the background fields undergo a $G_F$ gauge transformation
$A\to A^g$ and transform the coset representative by
$\Omega\to g\Omega$.  Applying (8.6.4) twice, or using (8.6.7), gives
\[
  \mathcal L_{\rm WZ}(A^g,g\Omega)-\mathcal L_{\rm WZ}(A,\Omega)
  =-\mathcal C_{\mathcal R}\,\alpha_4(A,g)
  +\dd(\cdots).
  \tag{8.6.8}
\]
The right-hand side is exactly the consistent anomaly obtained by varying the
fermion determinant, with the sign selected by whether the Wess--Zumino term
is added as inflow or as a compensating action.  Therefore the NGB functional
has the same anomalous gauge variation as the original $2M$ chiral fermions.

% BANKS-SOURCE: pdf=144 print=134 kind=solution id=8.7
\BanksSolution{8.7}
Choose the SU(5) generator used for weak hypercharge to be
\[
  Y=\operatorname{diag}\left(-\frac23,-\frac23,-\frac23,1,1\right),
  \qquad \operatorname{tr}Y=0 .
  \tag{8.7.1}
\]
The fundamental and anti-fundamental decompose as
\[
  5=(3,1,-2/3)\oplus(1,2,1),qquad
  \bar5=(\bar3,1,2/3)\oplus(1,2,-1).
  \tag{8.7.2}
\]
The 10 is the antisymmetric product of two fundamentals.  Its three types of
indices have charges
\[
  10=(\bar3,1,-4/3)\oplus(3,2,1/3)\oplus(1,1,2).
  \tag{8.7.3}
\]
The first term in (8.2) is the anti-down quark, and its second term is the
left-handed lepton doublet.  The three terms in (8.3) are respectively the
left-handed anti-up quark, the left-handed quark doublet, and the left-handed
anti-charged-lepton singlet.  These are precisely the fifteen left-handed
Weyl fields in one standard-model generation.  If a hypercharge is assigned
to the corresponding right-handed particle before charge conjugation, the sign
of the anti-up entry is reversed; (8.1)--(8.3) use the left-handed Weyl
convention relevant for anomaly sums.

As a consistency check, the cubic SU(5) anomaly coefficients obey
$A(10)=1$ and $A(\bar5)=-1$, so the reducible representation
$\bar5\oplus10$ is itself anomaly-free.

% BANKS-SOURCE: pdf=144 print=134 kind=solution id=8.8
\BanksSolution{8.8}
For SU($N$), normalize the fundamental anomaly to one.  The two-index
antisymmetric representation has
\[
  A(\Lambda^2N)=N-4,
  \qquad A(N)=1 .
  \tag{8.8.1}
\]
Thus the gauge anomaly cancels when
\[
  N-4+N_F=0,
  \qquad N_F=4-N,
  \tag{8.8.2}
\]
where the signed notation in the question treats a negative $N_F$ as
anti-fundamentals.  For the intended chiral construction $N\geq5$, put
$n=N-4$ and denote the
fundamental-flavor fields by $\chi^r$, $r=1,\ldots,n$; they transform as
anti-fundamentals of the gauge group.
The standard chiral construction is intended for $N\geq5$.  At $N=4$ the
antisymmetric representation is the real ${\bf6}$, so the SU(4)$^3$ gauge
anomaly already vanishes and $N_F=0$ is sufficient; its phase still has a
mixed SU(4)$^2$U(1) anomaly.  At $N=3$, $\Lambda^2{\bf3}\simeq\bar{\bf3}$,
and $N_F=1$ gives a vectorlike pair.  SU(2) has no perturbative cubic gauge
anomaly. Here $\Lambda^2{\bf2}$ is a gauge singlet, while the two doublets
suggested by $4-N=2$ give an even Witten-anomaly count. Pseudoreality also
enhances their flavor symmetry, so the $N=2$ theory must be treated separately
from the $U(1)_\psi\times U(n)_\chi$ analysis below.

For $N\geq5$, the kinetic action has the classical flavor group
\[
  U(1)_\psi\times U(n)_\chi
  =U(1)_\psi\times SU(n)_\chi\times U(1)_\chi .
  \tag{8.8.3}
\]
The SU($n$) factor is non-anomalous with respect to the SU($N$) gauge group,
because every SU($n$) generator is traceless.  For a general phase rotation
with charges $q_\psi$ and $q_\chi$, the mixed SU($N$)$^2$ anomaly is
\[
  \mathcal A_{SU(N)^2U(1)}
  =\frac12\bigl[(N-2)q_\psi+nq_\chi\bigr],
  \tag{8.8.4}
\]
where $T(\Lambda^2N)=(N-2)/2$ and $T(N)=1/2$.  There is one anomaly-free
continuous U(1), which can be chosen as
\[
  U(1)_X:\qquad q_\psi=n,qquad q_\chi=-(N-2).
  \tag{8.8.5}
\]
The independent U(1) combination is anomalous.  For example, the common
phase $q_\psi=q_\chi=1$ has coefficient $(N-2+n)/2$ in (8.8.4).  The
classical symmetry is therefore reduced quantum mechanically, with respect
to the gauge dynamics, to $SU(n)_\chi\times U(1)_X$ (apart from discrete
subgroups and the special case $n=0$).

The remaining triangle coefficients are 't Hooft anomalies of this global
group.  Set $\operatorname{tr}_{\boldsymbol n}(t^at^b)=\frac12\delta^{ab}$
and normalize the fundamental SU($n$) cubic anomaly to one.  Only $\chi$
carries SU($n$) flavor, with an additional multiplicity $N$ from its gauge
index.  Consequently,
\[
  \mathcal A_{SU(n)^3}=N,
  \qquad
  \mathcal A_{U(1)_XSU(n)^2}
    =-\frac{N(N-2)}2,
  \tag{8.8.6}
\]
while the coefficients with one SU($n$) generator and two U(1) insertions
vanish by tracelessness.  The cubic U(1) anomaly is
\begin{align*}
  \mathcal A_{U(1)_X^3}
    &=\frac{N(N-1)}2\,n^3+Nn\bigl[-(N-2)\bigr]^3\\
    &=-\frac{N^3(N-4)(N-3)}2
    \qquad (n=N-4).
\end{align*}
The signs of these global anomalies are reversed if every Weyl field is
replaced by its conjugate.  Equations (8.8.4)--(8.8.6) give the requested
three-current anomalies in a fixed normalization.

% BANKS-SOURCE: pdf=144 print=134 kind=solution id=8.9
\BanksSolution{8.9}
Let $q_{\bar5}$ and $q_{10}$ be the charges of the two SU(5) multiplets.
The mixed anomaly is proportional to their Dynkin indices,
\[
  \mathcal A_{U(1)SU(5)^2}
  =\frac12q_{\bar5}+\frac32q_{10}.
  \tag{8.9.1}
\]
It vanishes for $q_{\bar5}=-3$ and $q_{10}=1$.  The cubic anomaly of this
combination is
\[
  5(-3)^3+10(1)^3=-125 .
  \tag{8.9.2}
\]
An SU(5)-singlet Weyl fermion of charge $+5$ contributes $5^3=125$ and
cancels (8.9.2).  The resulting multiplets are conventionally written
\[
  10_{1}\oplus\bar5_{-3}\oplus1_{5} .
  \tag{8.9.3}
\]

Now use the ten Euclidean Clifford matrices.  From
\[
  a_i=\frac{\Gamma_{2i-1}+\ii\Gamma_{2i}}{2},
  \qquad
  a_i^\dagger=\frac{\Gamma_{2i-1}-\ii\Gamma_{2i}}{2},
\]
the Clifford relation gives
\begin{align*}
  \{a_i,a_j\}
   &=\frac14\{\Gamma_{2i-1}+\ii\Gamma_{2i},
              \Gamma_{2j-1}+\ii\Gamma_{2j}\}=0,\\
  \{a_i,a_j^\dagger\}
   &=\frac14\{\Gamma_{2i-1}+\ii\Gamma_{2i},
              \Gamma_{2j-1}-\ii\Gamma_{2j}\}=\delta_{ij}.
  \tag{8.9.4}
\end{align*}
The number operators $a_i^\dagger a_i$ therefore have eigenvalues zero or
one.  The 25 bilinears obey
\[
  [a_i^\dagger a_j,a_k^\dagger a_l]
   =\delta_{jk}a_i^\dagger a_l-\delta_{il}a_k^\dagger a_j,
  \tag{8.9.5}
\]
so their traceless part is $\mathfrak{su}(5)$ and their trace is the commuting
U(1) number operator $\mathcal N=\sum_i a_i^\dagger a_i$.

Choose a Fock vacuum satisfying $a_i|s\rangle=0$.  Every state is uniquely
of the form
\[
  a_{i_1}^\dagger\cdots a_{i_k}^\dagger|s\rangle,
  \qquad 0\leq k\leq5,
\]
with distinct ordered indices.  There are
\[
  \sum_{k=0}^{5}\binom{5}{k}=2^5=32
  \tag{8.9.6}
\]
such states, which proves the dimension of the Dirac spinor representation.
The product $\Gamma_*\propto\Gamma_1\cdots\Gamma_{10}$ anticommutes with
each individual $\Gamma_\mu$, since moving one $\Gamma_\mu$ through the
other nine produces a minus sign.  It consequently commutes with every
SO(10) generator $[\Gamma_\mu,\Gamma_\nu]$.

Each $\Gamma_\mu$ changes $k$ by one, whereas $\Gamma_*$ acts with a sign
$(-1)^k$ up to an overall convention.  The even and odd Fock sectors are
therefore invariant, and each has dimension
\[
  \sum_{\substack{k=0\\ k\ {\rm even}}}^{5}\binom{5}{k}
  =\sum_{\substack{k=0\\ k\ {\rm odd}}}^{5}\binom{5}{k}=16 .
  \tag{8.9.7}
\]
They are the two irreducible chiral SO(10) spinors.  With the U(1) generator
chosen as
\[
  X=5-2\mathcal N,
  \tag{8.9.8}
\]
the even sector decomposes as
\[
  k=0:\ 1_{5},\qquad k=2:\ 10_{1},\qquad k=4:\bar5_{-3}.
  \tag{8.9.9}
\]
Thus it is exactly (8.9.3): the standard-model generation plus the singlet.
The odd sector is the conjugate 16, with all charges reversed.  A Higgs VEV
in a suitable SO(10) representation can preserve SU(5)$\times$U(1), and then
the commuting generator is the $X$ in (8.9.8).

% BANKS-SOURCE: pdf=145 print=135 kind=solution id=8.10
\BanksSolution{8.10}
Only left-handed SU(2) doublets contribute to a mixed global
U(1)$SU(2)^2$ anomaly.  For one generation there are three colored quark
doublets, each with baryon number $1/3$, and one lepton doublet with lepton
number $1$.  With $T(2)=1/2$,
\[
  \mathcal A_{B\,SU(2)^2}
   =3\left(\frac13\right)\frac12=\frac12,
  \qquad
  \mathcal A_{L\,SU(2)^2}=1\cdot\frac12=\frac12 .
  \tag{8.10.1}
\]
For three generations both coefficients are multiplied by three.  Their
difference cancels because
\[
  \mathcal A_{(B-L)\,SU(2)^2}
  =\left[3\left(\frac13\right)-1\right]\frac12=0 .
  \tag{8.10.2}
\]
The same result says that an SU(2) instanton changes $B$ and $L$ equally,
while preserving $B-L$.

The commuting SO(10) generator in (8.9.9) is $X$, with charges
$X(10)=1$, $X(\bar5)=-3$, and $X(1)=5$.  In the standard left-handed
hypercharge convention, with the source's doubled hypercharge denoted by
$Y$, the physical charge is
\[
  B-L=\frac{X+2Y}{5} .
  \tag{8.10.3}
\]
For example, the charges of
$(q_L,u^c,e^c,d^c,l_L,\nu^c)$ are respectively
\[
  \frac13,\ -\frac13,\ 1,\ -\frac13,\ -1,\ 1,
\]
as (8.10.3) gives on the entries of $10_1\oplus\bar5_{-3}\oplus1_5$.
The hypercharge generator belongs to SU(5), so modulo an SU(5) generator the
anomaly-free B--L direction is represented by the extra commuting U(1), which
is the intended identification in the problem.

% BANKS-SOURCE: pdf=145 print=135 kind=solution id=8.11
\BanksSolution{8.11}
Let $p_1$ and $p_2$ be the incoming electron and positron momenta, and let
$k_1$ and $k_2$ be the outgoing $W^-$ and $W^+$ momenta.  Set
\[
\begin{aligned}
  q&=p_1+p_2=k_1+k_2,\qquad s=q^2,\qquad t=(p_1-k_1)^2,\\
  p_1^2&=p_2^2=0,\qquad k_1^2=k_2^2=-m_W^2,\qquad \widehat s=-s>0,\\
  P_{L,R}&=\frac12(1\mp\gamma_5).
\end{aligned}
\tag{8.11.1}
\]
We use $\eta=\operatorname{diag}(-,+,+,+)$ and
$p\cdot x=-p^0x^0+\mathbf p\cdot\mathbf x$, so the positive-frequency
factor is $\ee^{+\ii p\cdot x}$.  The $\gamma_5$ projectors in this
calculation are the $D=4$ chiral projectors.
For reference, with all three momenta incoming to the neutral-gauge-boson
vertex, use
\[
  \Gamma_{\mu\alpha\beta}(q,k_-,k_+)
  =\eta_{\alpha\beta}(k_--k_+)_\mu
   +\eta_{\beta\mu}(k_+-q)_\alpha
   +\eta_{\mu\alpha}(q-k_-)_\beta .
  \tag{8.11.2}
\]
The signs of every amplitude below change together if the opposite convention
for the interaction Lagrangian is used.

The left-handed electron has a $t$-channel neutrino graph.  Its invariant
amplitude is
\[
  \mathcal M_t^L
  =-\frac{g_2^2}{2}\,
  \bar v(p_2)\gamma^\beta P_L
  \frac{\slashed p_1-\slashed k_1}{t-\ii\epsilon}
  \gamma^\alpha P_Lu(p_1)
  \epsilon^*_{\alpha}(k_1)\epsilon^*_{\beta}(k_2).
  \tag{8.11.3}
\]
There is no such graph for a right-handed electron.  The two $s$-channel
graphs can be written together as
\[
  \mathcal M_s^h
  =\bar v(p_2)\gamma^\mu P_hu(p_1)
  \,\Gamma_{\mu\alpha\beta}(q,-k_1,-k_2)
  \,\epsilon^*_{\alpha}(k_1)\epsilon^*_{\beta}(k_2)
  \left[\frac{e^2Q_e}{s-\ii\epsilon}
  +\frac{g_{eh}^Zg_{WWZ}}{s+m_Z^2-\ii\epsilon}\right],
  \tag{8.11.4}
\]
where $Q_e=-1$, $g_{WWZ}=g_2\cos\theta_W$, and
\[
  g_{eL}^Z=\frac{g_2}{\cos\theta_W}
  \left(-\frac12+\sin^2\theta_W\right),
  \qquad
  g_{eR}^Z=\frac{g_2}{\cos\theta_W}\sin^2\theta_W .
  \tag{8.11.5}
\]
Thus $\mathcal M^L=\mathcal M_t^L+\mathcal M_s^L$, while
$\mathcal M^R=\mathcal M_s^R$.

In a general $R_\kappa$ gauge the neutral-vector propagator is
\[
  D_V^{\mu\nu}(q)=\frac{-\ii}{q^2+m_V^2-\ii\epsilon}
  \left[\eta^{\mu\nu}-(1-\kappa)
  \frac{q^\mu q^\nu}{q^2+\kappa m_V^2-\ii\epsilon}\right],
  \tag{8.11.6}
\]
with $m_\gamma=0$.  The current in (8.11.4) is conserved for a massless
electron:
\[
  q_\mu\bar v(p_2)\gamma^\mu P_hu(p_1)
  =\bar v(p_2)(\slashed p_1+\slashed p_2)P_hu(p_1)=0,
  \tag{8.11.7}
\]
because $\slashed p_1u(p_1)=0$ and $\bar v(p_2)\slashed p_2=0$, with the
projector moved through the first slash when needed.  Every $\kappa$-dependent
term in (8.11.6) is therefore zero.  The neutrino graph has no gauge-boson
propagator, and the electron Yukawa coupling to the would-be Goldstone is zero
in the same approximation.  Equations (8.11.3)--(8.11.7) prove the explicit
$\kappa$ independence of both helicity amplitudes.

For a longitudinal $W$,
\[
  \epsilon_L^\mu(k)=\frac{k^\mu}{m_W}+O\!\left(\frac{m_W}{E}\right).
  \tag{8.11.8}
\]
Individual terms in (8.3)--(8.4) then appear to grow as powers of
$E/m_W$.  Contracting the first factor of $k_1$ or $k_2$ into the neutrino
numerator and into (8.11.2), and using
$\slashed p_1u=\bar v\slashed p_2=0$, cancels the propagator terms between the
$t$ graph and the photon graph.  The remaining terms cancel against the $Z$
graph precisely because
\[
  e=g_2\sin\theta_W,\qquad g_{WWZ}=g_2\cos\theta_W,\qquad
  m_W=m_Z\cos\theta_W .
  \tag{8.11.9}
\]
This is the tree-level Ward identity of the broken gauge theory.  Equivalently,
the equivalence theorem gives
\[
  \mathcal M(e^+e^-\to W_L^+W_L^-)
  =\mathcal M(e^+e^-\to\phi^+\phi^-)
  +O\!\left(\frac{m_W}{E}\right),
  \tag{8.11.10}
\]
where $\phi^\pm$ are the charged would-be Goldstone bosons.  The scalar-pair
amplitude is of order $g_2^2$ at fixed angle and has no positive power of
$E/m_W$.  Hence the longitudinal production amplitude stays bounded at high
energy, and its cross section falls with the usual inverse power of $\widehat s$.

% BANKS-SOURCE: pdf=145 print=135 kind=solution id=8.12
\BanksSolution{8.12}
This determinant argument is four-dimensional and Euclidean, so the chiral
$\gamma_5$ pairing is used only in $D=4$.
Choose Euclidean gamma matrices Hermitian and a covariant derivative for which
$\slashed D$ is anti-Hermitian.  In the Weyl representation,
\[
  \slashed D=
  \begin{pmatrix}
    0&\sigma^\mu D_\mu\\
    \sigma_\mu^\dagger D_\mu&0
  \end{pmatrix} .
  \tag{8.12.1}
\]
The determinant of the Dirac operator with a positive real mass is, up to the
field-independent phase from the factors of $\ii$ in the question,
\[
  \det(\slashed D+m)
  =\det\!\left[
    m^2-(\sigma^\mu D_\mu)(\sigma_\nu^\dagger D_\nu)
  \right]
  =\det(m^2-\Delta_R).
  \tag{8.12.2}
\]
This follows from the Schur-complement identity applied to (8.12.1).  The
operator $\Delta_R$ is negative semidefinite, since
\[
  \langle\chi,\Delta_R\chi\rangle
  =-\|\sigma_\nu^\dagger D_\nu\chi\|^2\leq0 .
  \tag{8.12.3}
\]
Every nonzero eigenvalue of $\slashed D$ occurs in a pair
$+\ii\lambda$ and $-\ii\lambda$, related by $\gamma_5$.  Its contribution to
the determinant is therefore
\[
  (m+\ii\lambda)(m-\ii\lambda)=m^2+\lambda^2>0,
  \tag{8.12.4}
\]
and a zero mode contributes $m>0$.  Thus
\[
  \det(\ii\slashed D+\ii m)
  =\ii^{\dim\mathcal H}\det(\slashed D+m)
\]
is formally positive after removing the universal, field-independent phase
associated with the displayed convention.  A gauge-field-dependent sign or
phase is absent.

% BANKS-SOURCE: pdf=145 print=135 kind=solution id=8.13
\BanksSolution{8.13}
This static-potential calculation is four-dimensional.  Use the target
mostly-plus metric $\eta=\operatorname{diag}(-,+,+,+)$.
Let two static sources at $\mathbf 0$ and $\mathbf L$ carry generators
$t_1^a$ and $t_2^a$.  In this fixed $D=4$ static-source specialization,
their conserved currents are
\[
  J_1^{0a}(x)=g\,t_1^a\delta^3(\mathbf x),
  \qquad
  J_2^{0a}(x)=g\,t_2^a\delta^3(\mathbf x-\mathbf L),
  \qquad J_i^{ka}=0 .
  \tag{8.13.1}
\]
The leading covariant-gauge propagator is
\[
  D_{\mu\nu}{}^{ab}(q)=\frac{-\ii\delta^{ab}}{q^2-\ii\epsilon}
  \left[\eta_{\mu\nu}-(1-\kappa)\frac{q_\mu q_\nu}{q^2-\ii\epsilon}\right].
  \tag{8.13.2}
\]
The static currents impose $q^0=0$.  The gauge-dependent term then vanishes
already because $q_0=0$; more generally it vanishes after contraction with any
conserved source, $q_\mu J^\mu=0$.  The cross term in the source action, in
this fixed $D=4$ specialization, is
\[
  V_{\rm int}(L)=g^2t_1^at_2^a
  \int\frac{\dd^3\mathbf q}{(2\pi)^3}
  \frac{\ee^{\ii\mathbf q\cdot\mathbf L}}{\mathbf q^2}
  =\frac{g^2}{4\pi L}t_1^at_2^a .
  \tag{8.13.3}
\]
The two same-line contractions produce divergent self-energies independent of
$L$.  Non-abelian three- and four-gauge vertices first affect this result at
order $g^4$.  Consequently the leading static potential is
\[
  V(L)=V_{\rm self}+\frac{g^2}{4\pi L}t_1^at_2^a+O(g^4),
  \tag{8.13.4}
\]
independent of $\kappa$.  For a source and anti-source in $R$ and $\bar R$,
$t_2^a=t_{\bar R}^a$, and (8.13.4) agrees with the Wilson-loop result.

% BANKS-SOURCE: pdf=145 print=135 kind=solution id=8.14
\BanksSolution{8.14}
This equal-time current algebra is the $D=4$ chiral calculation.
Use equal-time canonical anticommutation relations, for example
\[
  \{e_\alpha(t,\mathbf x),e_\beta^\dagger(t,\mathbf y)\}
  =\delta_{\alpha\beta}\delta^3(\mathbf x-\mathbf y),
  \tag{8.14.1}
\]
with analogous relations for each species and zero anticommutators between
different species.  Put $A=1-\gamma_5$.  For one lepton doublet,
\[
  j_+^0=e^\dagger A\nu_e,\qquad j_-^0=\nu_e^\dagger A e,
\]
and the bilinear identity following from (8.14.1) is
\[
  [e^\dagger A\nu_e,\nu_e^\dagger A e]
  =\delta^3(\mathbf x-\mathbf y)
  \left(e^\dagger A^2e-\nu_e^\dagger A^2\nu_e\right).
  \tag{8.14.2}
\]
Since $A^2=2A$, adding the electron and muon doublets gives, in this fixed
$D=4$ specialization,
\[
\begin{split}
  [J_+^0(t,\mathbf x),J_-^0(t,\mathbf y)]
  ={}&2\delta^3(\mathbf x-\mathbf y)\bigl[
  e^\dagger(1-\gamma_5)e-\nu_e^\dagger(1-\gamma_5)\nu_e\\
  &\hspace{6.6em}
  +\mu^\dagger(1-\gamma_5)\mu
  -\nu_\mu^\dagger(1-\gamma_5)\nu_\mu\bigr](t,\mathbf x).
\end{split}
  \tag{8.14.3}
\]
In terms of left-handed fields this is proportional to
$e_L^\dagger e_L-\nu_{eL}^\dagger\nu_{eL}\,+\mu_L^\dagger\mu_L-
\nu_{\mu L}^\dagger\nu_{\mu L}$, the third generator of
the weak-isospin algebra.  The electromagnetic density instead contains the
electric charges of both chiralities of the charged leptons and no neutrino
density:
\[
  J_{\rm em}^0=-e^\dagger e-\mu^\dagger\mu+\cdots .
  \tag{8.14.4}
\]
Thus the charge-changing currents do not close with the electromagnetic
current.  An independent commuting U(1) generator is required, giving the
SU(2)$\times$U(1) electroweak gauge group.

% BANKS-SOURCE: pdf=146 print=136 kind=solution id=8.15
\BanksSolution{8.15}
The Ward identities and the effective functional below are written for general
spacetime dimension $D$.
Let $[T^a,T^b]=\ii f^{ab}{}_cT^c$, and let $J_a{}^\mu$ be the corresponding
currents.  A local change of variables in a correlation function gives the
Ward identities
\begin{equation}
\begin{aligned}
  \partial_\mu^x\left\langle J_a{}^\mu(x)
  \prod_{r=1}^nJ_{a_r}{}^{\mu_r}(x_r)\right\rangle
  ={}&-\ii\sum_{r=1}^n f_{a a_r}{}^c\delta^D(x-x_r)\\
  &\times\left\langle J_{a_1}{}^{\mu_1}(x_1)\cdots
  J_c{}^{\mu_r}(x_r)\cdots J_{a_n}{}^{\mu_n}(x_n)\right\rangle .
\end{aligned}
\tag{8.15.1}
\end{equation}
with the usual contact-term convention for the insertion at $x_r$.  Define
\[
  Z[A]=\left\langle\exp\!\left(\ii\int\dd^D x\,
  J_a{}^\mu A_\mu{}^a\right)\right\rangle=\ee^{\ii W[A]}.
\]
Summing (8.15.1) gives the functional identity
\[
  D_\mu{}^{ab}(A)\frac{\delta W}{\delta A_\mu{}^b(x)}=0,
  \qquad
  D_\mu{}^{ab}(A)=\delta^{ab}\partial_\mu+f^{acb}A_\mu{}^c,
  \tag{8.15.2}
\]
up to the corresponding convention for the factors of $\ii$.  Conversely,
integrating (8.15.2) against an arbitrary compactly supported
$\omega^a(x)$ gives
\[
  \delta_\omega W[A]=\int\dd^D x\,
  \frac{\delta W}{\delta A_\mu{}^a}D_\mu{}^{ab}(A)\omega^b=0,
  \tag{8.15.3}
\]
so the Ward identities and gauge invariance of $W$ are equivalent.

The connected-current expansion is the cumulant expansion
\[
\begin{aligned}
  W[A]={}&\sum_{n\geq1}\frac{\ii^{\,n-1}}{n!}
  \int\!\dd^D x_1\cdots\dd^D x_n\,
  A_{\mu_1}{}^{a_1}(x_1)\cdots A_{\mu_n}{}^{a_n}(x_n)\\
  &\times\left\langle J_{a_1}{}^{\mu_1}(x_1)\cdots
  J_{a_n}{}^{\mu_n}(x_n)\right\rangle_{\!c} .
\end{aligned}
\tag{8.15.4}
\]
In particular,
\[
  \left.\frac{\delta^nW}{\delta A_1\cdots\delta A_n}\right|_{A=0}
  =\ii^{\,n-1}\langle J_1\cdots J_n\rangle_c.
  \tag{8.15.5}
\]

For the non-standard transform, write the quadratic term with its spacetime
integral and invariant index contraction understood:
\[
  \Gamma[a]+M^2(A-a)^2=W[A].
  \tag{8.15.6}
\]
The stationary value of $A$ is determined by
\[
  \frac{\delta W}{\delta A_\mu{}^a(x)}
  =2M^2\bigl(A_\mu{}^a(x)-a_\mu{}^a(x)\bigr),
  \qquad
  \frac{\delta\Gamma}{\delta a_\mu{}^a(x)}
  =2M^2\bigl(A_\mu{}^a(x)-a_\mu{}^a(x)\bigr).
  \tag{8.15.7}
\]
To see the relation to connected functions explicitly, let
$W=\frac12A K A+O(A^3)$, where $K$ is the connected two-current kernel in
the convention of (8.15.4).  Then
\[
  A=\frac{2M^2}{2M^2-K}\,a+O(a^2),
  \qquad
  \Gamma^{(2)}=\frac{2M^2K}{2M^2-K}.
  \tag{8.15.8}
\]
At cubic order, for example,
\[
  \Gamma^{(3)}=W^{(3)}
  \left(\frac{2M^2}{2M^2-K}\right)^3,
  \tag{8.15.9}
\]
with all indices and momenta understood.  At fourth and higher order there
are additional terms made from products of lower connected kernels joined by
$ (2M^2-K)^{-1}$.  Thus the coefficients of $\Gamma$ are the connected
current vertices with external legs dressed by the quadratic transform and
with one-particle-reducible chains removed, the usual effective-action
organization.

If $A$ and $a$ are transformed by the same gauge transformation, the quadratic
term in (8.15.6) is invariant.  The stationary equation (8.15.7) maps a
solution $A[a]$ into the solution $A^g[a^g]$, and the Ward identity (8.15.2)
therefore implies
\[
  \Gamma[a^g]=\Gamma[a].
  \tag{8.15.10}
\]
At momenta below the mass of any omitted state, locality and gauge invariance
give the Yang--Mills form at leading order,
\[
  \Gamma[a]=-\frac1{4g_{\rm eff}^2}\int\dd^D x\,
  F_{\mu\nu}{}^a(a)F^{a\mu\nu}+\text{higher-dimensional terms},
  \tag{8.15.11}
\]
with a possible topological $\theta$ term in the $D=4$ specialization.

Suppose now that $G$ is spontaneously broken to $H$.  Let $\Omega(x)$ be a
coset representative, with the redundancy $\Omega\sim\Omega h$ and with a
left gauge transformation $\Omega\to g\Omega$.  The gauge-covariant Maurer--
Cartan field is
\[
  B_\mu=\Omega^{-1}(\partial_\mu-gA_\mu)\Omega,
  \qquad B_\mu=a_\mu+W_\mu,
  \tag{8.15.12}
\]
where $a_\mu$ lies in $\mathfrak h$ and $W_\mu$ in the coset subspace.  Under
$\Omega\to\Omega h$, $a_\mu$ is an $H$ connection and $W_\mu$ transforms
homogeneously.  Consequently the lowest-derivative coupling of the NGBs to
the vector field is
\[
  \mathcal L_{\rm NGB}
  =-\frac{f^2}{2}\,\langle W_\mu,W^\mu\rangle
  =-\frac{f^2}{2}\,\operatorname{tr}_{\mathfrak g/\mathfrak h}
  \bigl[(\partial_\mu\Omega-gA_\mu\Omega)^\dagger
  (\partial^\mu\Omega-gA^\mu\Omega)\bigr],
  \tag{8.15.13}
\]
where the second notation means the invariant projection onto the coset
directions.  Together, (8.15.11), (8.15.13), and the quadratic term in
(8.15.6) give the Yang--Mills theory of vector mesons interacting with the
massless NGBs.

In unitary gauge, $\Omega=1$, the NGB term becomes a vector mass term only in
the broken directions,
\[
  \mathcal L_{\rm NGB}\big|_{\Omega=1}
  =-\frac12\,a_\mu^a(\Delta m^2)_{ab}a^{b\mu},
  \qquad
  (\Delta m^2)_{ab}=g_ag_b f^2\,
  \langle T_a^\perp,T_b^\perp\rangle .
  \tag{8.15.14}
\]
Here $T_a^\perp$ is the projection of $T_a$ onto
$\mathfrak g/\mathfrak h$, and the invariant inner product is fixed by the
normalization of the NGB kinetic term.  If the quadratic transform contributes
the matrix $M^2_{ab}$, the physical vector mass matrix is therefore
\[
  (m_V^2)_{ab}=M^2_{ab}+(\Delta m^2)_{ab}
  =M^2_{ab}+g_ag_b f^2
  \langle T_a^\perp,T_b^\perp\rangle .
  \tag{8.15.15}
\]
The $H$ components have no term in (8.15.14); the broken components acquire
the additional mass supplied by the NGB decay constant.  This is the requested
contribution beyond the matrix $M^2$.

For the two-flavor chiral example,

\[
  G=SU(2)_L\times SU(2)_R,\qquad H=SU(2)_V,
  \qquad U(x)=\exp\!\left(\frac{2\ii\pi^a(x)\tau^a}{F_\pi}\right),
  \tag{8.15.16}
\]

the vector generators $T_V^a=T_L^a+T_R^a$ span $\mathfrak h$, while
$T_A^a=T_L^a-T_R^a$ span the coset.  The corresponding vector fields are
identified with the $\rho^a$ and $A_1{}^a$ multiplets.  The singlet vector
extension of the flavor group supplies the $\omega$.  Equation (8.15.13)
becomes the usual covariant pion kinetic term, where
$D_\mu U=\partial_\mu U-g_LA_{L\mu}U+g_RU A_{R\mu}$,

\[
  \mathcal L_\pi=-\frac{F_\pi^2}{4}
  \operatorname{tr}\bigl[(D_\mu U)^\dagger D^\mu U\bigr],
  \tag{8.15.17}
\]

and its expansion supplies the pion couplings to $\rho$, $\omega$, and $A_1$.
The $H$-vector masses come from $M^2$, whereas the axial mass block receives
the additional $g_A^2F_\pi^2$ contribution in (8.15.14).  This construction is
a vector-dominance model.  It has no controlled limit in which the vector
masses are parametrically below the chiral scale $4\pi F_\pi$, so replacing
the full $\Gamma$ by its Yang--Mills part is a phenomenological approximation;
the chiral Lagrangian itself has a low-momentum expansion with a controlled
small parameter.

\end{bankssolutions}
